%% ============================================================
%%  Chapitre 3 — Moteur de Backtest par Walk-Forward Analysis
%% ============================================================
\chapter{Moteur de Backtest par Walk-Forward Analysis}
\label{ch:wfo}

\section{Pourquoi la Walk-Forward Analysis~?}
\label{sec:pourquoi_wfa}

\subsection{Le problème central du backtesting classique}

Un backtest classique optimise les paramètres d'une stratégie sur un historique,
puis mesure la performance sur ce même historique. Cette approche est circulaire~:
les paramètres sont sélectionnés \emph{parce qu'ils fonctionnent} sur ces données,
non pas parce qu'ils sont genuinement robustes.

Pardo~\cite{pardo2008} pose le problème sans détour~:

\begin{quote}
\textit{«~An in-sample backtest is not a backtest — it is a fitting exercise.~»}
— Pardo (2008), Ch.~11.
\end{quote}

\subsection{La WFA comme simulation réaliste}

La Walk-Forward Analysis (WFA) résout ce problème en simulant la manière dont une
stratégie est réellement utilisée en production~:

\begin{quote}
\textit{«~Walk-Forward Analysis is the closest possible simulation of the way in
which an optimizable trading strategy is typically used in real time.~»}
— Pardo (2008), Ch.~11, p.~237.
\end{quote}

Elle répond à quatre questions essentielles~\cite{pardo2008}~:
\begin{enumerate}
    \item La stratégie est-elle robuste~? Fonctionnera-t-elle en trading réel~?
    \item Quel taux de rendement produira-t-elle en conditions réelles~?
    \item Comment les changements de comportement de marché affecteront-ils
    la performance~?
    \item Quels sont les meilleurs paramètres à utiliser actuellement~?
\end{enumerate}

\subsection{Motivation pratique — Semaine 3}

Les premiers tests effectués en Semaine~3 du stage~\cite{rapport_s3} ont illustré
concrètement ce problème~: une SMA optimisée sur 20 ans de données IAM avec de très
petites fenêtres produisait des milliers de transactions et des résultats d'apparence
excellente. Ces résultats étaient non réalistes~: la stratégie avait mémorisé le
passé plutôt qu'appris un comportement transférable.

\medskip

Par la suite, l'hypothèse d'un \textbf{régime haussier depuis 2023} sur le MASI~\cite{rapport_s4}
a orienté la fenêtre de calibration vers des périodes récentes, plus représentatives
du régime actuel. La WFA formalise et systématise cette intuition.

%% ============================================================
\section{La Fonction Objectif PROM}
\label{sec:prom}

\subsection{Pourquoi pas le net profit~?}

La métrique la plus naturelle pour évaluer une stratégie est son profit net. Pardo
démontre cependant que cette métrique est une cible inadéquate pour l'optimisation~:

\begin{quote}
\textit{«~Research has shown that highest net profit alone is not an adequate measure
of the overall trading performance and robustness of a trading model. For example,
60 percent of the highest net profit could be generated by a strategy with one large
and likely unrepeatable trade, such as a short trade during the October 1987 stock
market crash.~»} — Pardo (2008), Ch.~9, p.~202.
\end{quote}

\subsection{Définition formelle du PROM}

Le \textbf{PROM} (Pessimistic Return on Margin) est une variante pessimiste du
retour sur marge qui pénalise les stratégies à forte variance de trade~:

\begin{equation}
\text{PROM} = \frac{(W \cdot \bar{w}) - (L \cdot \bar{l})}{M}
\cdot \left(1 - \text{CF}\right)
\label{eq:prom}
\end{equation}

où~:
\begin{itemize}
    \item $W$~: nombre de trades gagnants~;
    \item $L$~: nombre de trades perdants~;
    \item $\bar{w}$, $\bar{l}$~: gain/perte moyens~;
    \item $M$~: marge maximale utilisée~;
    \item $\text{CF}$~: coefficient de friction (facteur de pénalité croissant avec
    la variance du PnL, favorisant la stabilité).
\end{itemize}

Le PROM sélectionne les configurations qui produisent des gains \emph{stables et
répétables}, plutôt que celles dont la performance est concentrée sur un ou deux
trades exceptionnels.

\begin{table}[H]
\centering
\caption{Comparaison des fonctions objectif sur un exemple MASI (données illustratives)}
\label{tab:objectif_compare}
\begin{tabularx}{\linewidth}{@{}lYYY@{}}
\toprule
\textbf{Critère} & \textbf{Net Profit} & \textbf{Sharpe Ratio} & \textbf{PROM} \\
\midrule
Stratégie A (strat. stable)   & 12\,\%  & 0.85 & 0.72 \\
Stratégie B (1 trade géant)   & 18\,\%  & 0.41 & 0.28 \\
\textbf{Gagnante selon critère} & B & A & \textbf{A} \\
\bottomrule
\multicolumn{4}{l}{\small La stratégie B a un net profit supérieur, mais sa performance dépend d'un trade exceptionnel.} \\
\multicolumn{4}{l}{\small PROM et Sharpe éliminent cet artefact ; PROM est la mesure officielle retenue.} \\
\end{tabularx}
\end{table}

%% ============================================================
\section{Le Pipeline WFO Complet}
\label{sec:pipeline_wfo}

\subsection{Vue d'ensemble}

La WFA est une procédure d'optimisation-validation en fenêtres glissantes. Pour
chaque titre dans la stratégie~:

\begin{enumerate}
    \item Diviser l'historique en fenêtres IS/OOS chevauchantes~;
    \item Dans chaque fenêtre IS~: optimiser les paramètres par PROM~;
    \item Dans chaque fenêtre OOS~: tester le gagnant IS sur données inédites~;
    \item Calculer le Walk-Forward Efficiency (WFE) sur toutes les fenêtres~;
    \item Accepter ou rejeter la stratégie selon WFE $\geq 50\,\%$~;
    \item Confirmer sur la période test finale (\textit{hold-out}).
\end{enumerate}

\begin{figure}[H]
\centering
\placeholder{Schéma détaillé du pipeline WFO :\\
Axe temporel horizontal : données historiques 2015–2025 + période test 2025+. \\
Zone WFO (2015–2024) : 6 fenêtres IS/OOS glissantes (IS en bleu, OOS en orange). \\
Zone test finale (2024–2025) : en rouge, jamais touchée pendant l'optimisation. \\
Flèche : "Résultats OOS agrégés → calcul WFE". \\
Flèche : "Paramètres du gagnant IS final → déploiement sur période test". \\
Légende claire.}
\caption{Pipeline complet de la Walk-Forward Analysis.
La période test finale (en rouge) est exclue de tout processus d'optimisation.}
\label{fig:wfo_pipeline}
\end{figure}

\subsection{Step 1 --- Génération des Variantes}

Les paramètres optimisables couvrent l'ensemble des composantes de la stratégie~:

\begin{itemize}
    \item \textbf{Paramètres d'indicateurs}~: fenêtres SMA/EMA/RSI, periods MACD, etc.~;
    \item \textbf{Seuils d'entrée/sortie}~: niveaux RSI overbought/oversold, etc.~;
    \item \textbf{Sizing}~: \texttt{size\_pct}, \texttt{reduction\_pct}, Kelly modifier~;
    \item \textbf{Risque}~: multiplicateur ATR stop, ratio récompense/risque~;
    \item \textbf{Timing}~: cooldown bars, time stop bars.
\end{itemize}

Pardo préconise des pas proportionnels à la taille du range de recherche~:

\begin{quote}
\textit{«~Excessively fine scanning of a parameter range can be a significant
contributing factor to overfitting.~»} — Pardo (2008), Ch.~10, p.~252.
\end{quote}

\begin{table}[H]
\centering
\caption{Calibration des pas de scan selon la taille du range (Pardo, Ch.~10)}
\label{tab:scan_steps}
\begin{tabularx}{\linewidth}{@{}YYY@{}}
\toprule
\textbf{Taille du range} & \textbf{Pas recommandé} & \textbf{Nb de candidats} \\
\midrule
5–40 (35 valeurs)    & 1 & 36 \\
20–100 (80 valeurs)  & 2 & 41 \\
50–250 (200 valeurs) & 5 & 41 \\
\bottomrule
\end{tabularx}
\end{table}

\begin{table}[H]
\centering
\caption{Plages de paramètres par horizon temporal}
\label{tab:params_horizons}
\begin{tabularx}{\linewidth}{@{}lYYY@{}}
\toprule
\textbf{Paramètre} & \textbf{Court terme} & \textbf{Moyen terme} & \textbf{Long terme} \\
\midrule
Fenêtre SMA/EMA             & 5–40   & 20–100  & 50–250 \\
ATR stop multiplier         & 1–3    & 1.5–4   & 2–5    \\
Reward-to-risk ratio        & 1–2.5  & 1.5–3.5 & 2–5    \\
Cooldown (barres)           & 1–5    & 3–15    & 5–30   \\
Time stop (barres)          & 5–20   & 15–60   & 40–120 \\
Seuils RSI (survendu)       & 25–40  & 25–40   & 25–40  \\
\bottomrule
\end{tabularx}
\end{table}

\subsection{Steps 2 et 3 --- Optimisation IS et Neighbor Averaging}

Dans chaque fenêtre IS, l'optimiseur recherche les paramètres maximisant le PROM.
Pour éviter de sélectionner un optimum \emph{isolé} dans l'espace des paramètres
(souvent signe d'overfitting), un \textbf{neighbor averaging} est appliqué~:
la performance d'une configuration est lissée avec ses voisins immédiats dans
l'espace de paramètres. Les paramètres à plateau large et stable sont préférés.

\subsection{Steps 4 à 6 --- Rolling OOS et Agrégation}

Le gagnant IS est appliqué mécaniquement sur la fenêtre OOS correspondante.
Les résultats OOS sont agrégés sur toutes les fenêtres pour produire une
distribution de performances OOS représentative de la robustesse.

\subsection{Step 7 --- Critère d'Acceptation WFE}

Le Walk-Forward Efficiency est défini comme~:

\begin{equation}
\text{WFE} = \frac{\text{Retour annualisé OOS}}{\text{Retour annualisé IS}}
\label{eq:wfe}
\end{equation}

La règle d'acceptation~:

\begin{equation}
\text{Stratégie acceptée} \iff \text{WFE} \geq 50\,\%
\label{eq:wfe_critere}
\end{equation}

Pardo justifie ce seuil~:

\begin{quote}
\textit{«~Research has clearly demonstrated that robust trading strategies have WFEs
greater than 50 or 60 percent and in the case of extremely robust strategies, even
higher.~»} — Pardo (2008), Ch.~11, p.~239.
\end{quote}

\figuretodo{wfe_barplot}{Barplot du WFE par action pour un exemple de portefeuille MASI.
Ligne horizontale rouge à WFE $= 50\,\%$. Actions acceptées (vert) vs rejetées (rouge).}

\subsection{Steps 8 et 9 --- Paramètres Finaux et Période Test}

La dernière fenêtre IS fournit les paramètres finaux à déployer. Ces paramètres
sont ensuite validés sur la \textbf{période test hold-out}~: des données qui n'ont
jamais été impliquées dans aucune des fenêtres d'optimisation.

%% ============================================================
\section{Validation Statistique}
\label{sec:validation_stat}

\subsection{Deflated Sharpe Ratio}

Le Sharpe Ratio classique ne tient pas compte du nombre de stratégies testées.
Bailey \& López de Prado~\cite{bailey2012} introduisent le Deflated Sharpe Ratio
(DSR), qui ajuste le Sharpe pour le nombre de trials~:

\begin{equation}
\text{DSR}(\hat{SR}) = \Phi\!\left(
\frac{(\hat{SR} - SR^*) \sqrt{T-1}}{\sqrt{1 - \hat{\gamma}_3 \hat{SR}
+ \frac{\hat{\gamma}_4 - 1}{4}\hat{SR}^2}}
\right)
\label{eq:dsr}
\end{equation}

où $\Phi$ est la CDF de la loi normale, $SR^*$ est le Sharpe espéré du meilleur
parmi $N$ tentatives aléatoires, $\hat{\gamma}_3$ et $\hat{\gamma}_4$ sont
respectivement le skewness et l'excès de kurtosis des retours.

Le DSR est calculé sur les résultats OOS de la WFA et constitue la mesure de
significativité statistique de la stratégie.

\subsection{Validation par Monte Carlo}

Pour compléter le DSR, une simulation Monte Carlo bootstrap est conduite~: les
retours OOS des fenêtres sont rééchantillonnés $10\,000$ fois pour construire
la distribution des performances sous l'hypothèse nulle. La p-value empirique
est calculée comme la fraction de simulations qui surpasse le Sharpe OOS observé.

\figuretodo{montecarlo_dist}{Distribution des Sharpe ratios Monte Carlo (10\,000 simulations bootstrap) pour la stratégie gagnante sur IAM. Ligne rouge : Sharpe OOS observé. Zone grise : zone de non-significativité. p-value affichée.}

\subsection{Kelly Sizing depuis les Trades OOS}

Le dimensionnement des positions est calibré non pas sur les trades IS (biaisés),
mais sur les \textbf{trades OOS} des fenêtres de Walk-Forward. La fraction de Kelly
partielle est appliquée~:

\begin{equation}
f^* = \frac{p \cdot b - q}{b}
\qquad \Rightarrow \qquad f_{\text{partiel}} = \kappa \cdot f^*
\label{eq:kelly}
\end{equation}

où $p$ est la probabilité de gain OOS, $b = \bar{w}/\bar{l}$ est le ratio
gain moyen/perte moyenne OOS, $q = 1-p$, et $\kappa \in (0,1]$ est le
Kelly modifier (paramètre configurable, typiquement 0.5 pour une fraction
demi-Kelly).

%% ============================================================
\section{Optimisations Techniques}
\label{sec:optim_perf}

\subsection{Le goulot d'étranglement}

Une optimisation WFO typique sur 93 tickers avec 1000 trials par fenêtre et
6 fenêtres requiert $93 \times 1\,000 \times 6 = 558\,000$ évaluations de
stratégie. Sans optimisation, ce calcul prenait plusieurs dizaines de minutes.

\subsection{Kernel Numba JIT}

Le goulot d'étranglement a été isolé dans la boucle de calcul de statistiques
par trial. Un kernel Numba JIT compilé à la volée (\texttt{@njit(cache=True)})
remplace la boucle Python~:

\begin{lstlisting}[caption={Kernel Numba JIT pour le calcul de métriques par trial},label={lst:numba_kernel}]
@njit(cache=True)
def _run_stats_fast_single_nb(
    signal_arr: np.ndarray,
    returns_arr: np.ndarray,
    cost_per_trade: float,
) -> tuple[float, float, float, float]:
    """
    Retourne (sharpe, max_drawdown, win_rate, prom) pour un trial.
    Aucun remplissage de positions (fills) n'est calculé ici.
    """
    # ... calcul vectorisé des statistiques uniquement
    return sharpe, max_drawdown, win_rate, prom
\end{lstlisting}

\textbf{Principe}~: seules 4 scalaires sont calculées par trial (pas de
tableaux de positions, pas de rapport). Le rapport complet n'est généré
qu'une seule fois pour le \emph{gagnant}.

\subsection{Hoisting de l'Indicateur Bank}

Les indicateurs sont précalculés une seule fois avant la boucle d'optimisation,
et non recalculés à chaque trial~:

\begin{lstlisting}[caption={Hoisting de l'indicateur bank avant la boucle de trials},label={lst:hoisting}]
# Calcul une seule fois (avant la boucle de trials)
indicator_bank = build_bank(ohlcv, strategy_config)

# Boucle d'optimisation : ~420 trials/s
for params in iter_trials():
    signal = make_signal_fast(indicator_bank, params)   # numpy, rapide
    sharpe, mdd, wr, prom = _run_stats_fast_single_nb(signal, returns, cost)
    if prom > best_prom:
        best_params = params
        best_prom = prom

# Un seul appel complet pour le gagnant
report = BacktestEngine.run(best_params, ohlcv, ...)
\end{lstlisting}

\subsection{Résultats de l'Optimisation}

\begin{table}[H]
\centering
\caption{Benchmarks avant/après optimisation technique (IAM, 1000 trials, machine de référence)}
\label{tab:benchmarks}
\begin{tabularx}{\linewidth}{@{}lYYY@{}}
\toprule
\textbf{Métrique}      & \textbf{Avant} & \textbf{Après} & \textbf{Gain} \\
\midrule
Latence cold JIT (1er appel) & 1\,808 ms & 47.6 ms  & $\mathbf{38\times}$ \\
Latence warm cache           & 288 ms    & 47.6 ms  & $6.1\times$          \\
Trials / seconde             & 69.5      & 420.5    & $6.1\times$          \\
Tests unitaires passant       & 73 / 73   & 73 / 73  & — \\
\bottomrule
\end{tabularx}
\end{table}

Ces gains permettent d'exécuter une optimisation WFO complète sur un ticker en
quelques secondes, rendant le traitement batch de l'ensemble du MASI pratique
en conditions réelles.

%% ============================================================
\section*{Conclusion du chapitre}

Ce chapitre a présenté le second moteur de résultat de la plateforme~: le moteur
de Backtest WFO. En implémentant fidèlement la Walk-Forward Analysis de Pardo avec
PROM comme fonction objectif, et en ajoutant une validation statistique par DSR et
Monte Carlo, ce moteur fournit un verdict de robustesse quantifié sur chaque stratégie.
L'optimisation technique par Numba JIT rend ce processus viable en temps de calcul
industriel. Le chapitre suivant présente l'architecture technique qui héberge et
orchestre ces deux moteurs.
